\section{Background}
\label{section:background}

This chapter establishes the technical foundations on which the rest of the thesis rests. Section~\ref{sec:bg:lvlm} describes how large vision-language models are constructed and how they generate text conditioned on an image. Section~\ref{sec:bg:medical-lvlm} covers their adaptation to medicine and, importantly, the failure modes that domain pretraining does not remove. Section~\ref{sec:bg:alignment} develops the alignment pipeline from supervised fine-tuning through reinforcement learning from human feedback to Direct Preference Optimization, deriving the DPO objective in full because the structure of that derivation is what the analysis in Chapter~\ref{section:analysis} exploits. Section~\ref{sec:bg:dpo-variants} surveys the variants that address DPO's known weaknesses. Section~\ref{sec:bg:kl} develops the information-theoretic machinery needed to reason about the regularizer, and Section~\ref{sec:bg:benchmarks} describes the datasets and metrics used to evaluate clinical vision-language systems.

\subsection{Large Vision-Language Models}
\label{sec:bg:lvlm}

A large vision-language model (LVLM) is a system that accepts an interleaved sequence of images and text and produces text. Contemporary LVLMs are overwhelmingly built by composing three components that are pretrained separately and then connected: a visual encoder, a modality projector, and an autoregressive language model. Understanding this decomposition matters for this thesis because the alignment method developed in Chapter~\ref{section:method} modifies only the third component, and the interpretation of its grounding results depends on that fact.

\subsubsection{Visual Encoders and Modality Projection}

The visual encoder maps an input image to a sequence of continuous vectors. The dominant choice is a Vision Transformer (ViT)~\cite{dosovitskiy2021image}, which partitions the image into fixed-size patches, embeds each patch linearly, adds positional information, and processes the resulting sequence with the same self-attention architecture used for text~\cite{vaswani2017attention}. The output is one vector per patch, so an image becomes a sequence of \emph{visual tokens} whose length is determined by the image resolution and the patch size.

The quantities involved are worth making concrete, because they explain several practical constraints encountered later. At a patch size of $14$ or $16$ pixels, a $1024 \times 1024$ image yields on the order of four thousand patches before any merging. Architectures typically reduce this --- for instance by merging adjacent $2 \times 2$ patch groups into a single token --- but a single high-resolution medical image still occupies hundreds to low thousands of positions in the language model's context. Since attention cost grows quadratically in sequence length, and since the objective developed in this thesis requires several forward passes over image-bearing sequences per optimization step, the visual token count is the dominant term in both the memory and the time budget reported in Section~\ref{sec:met:implementation}.

Encoders are typically initialized from contrastive image-text pretraining, most influentially CLIP~\cite{radford2021learning}, which trains an image encoder and a text encoder jointly so that matched image-caption pairs receive high cosine similarity and mismatched pairs low similarity, over very large web corpora. The resulting visual representation is already aligned to language in a coarse, whole-image sense, which is why it transfers well as the perceptual front end of a generative model.

It is worth noting for later that this pretraining objective supervises \emph{global} correspondence between an image and a caption. The loss is computed on a single pooled embedding per image; nothing in it requires that individual regions be separable, or that a small region be represented distinctly from the background it sits in. This is a reasonable trade for natural images, where the caption usually describes the dominant object. It is a poor fit for medical imaging, where the diagnostically decisive feature may occupy a fraction of a percent of the frame and differ from healthy tissue in texture rather than shape. The difficulty that medical LVLMs have with fine-grained localization is therefore not purely a matter of insufficient medical data; it is partly inherited from the pretraining objective of the encoder itself.

Because the visual encoder and the language model were trained independently, their representation spaces do not coincide. A \emph{projector} bridges them. The simplest and still widely used design is a linear layer or a small multilayer perceptron mapping each visual token into the language model's embedding dimension, as in LLaVA~\cite{liu2023visual, liu2024improved}. More elaborate designs compress the visual sequence before injection: BLIP-2's Q-Former~\cite{li2023blip2} uses a fixed set of learned query vectors that cross-attend to the encoder output, producing a short, fixed-length summary regardless of image size, and Flamingo~\cite{alayrac2022flamingo} instead interleaves gated cross-attention layers throughout a frozen language model.

These designs trade capacity against cost in a way that matters for grounding. A linear projector preserves the spatial correspondence between visual tokens and image regions: token $i$ still derives from patch $i$, so attention paid to that token is attention paid to a known location. A learned-query bottleneck destroys this correspondence, since each output vector is a weighted mixture over the whole image. The attention-based grounding measurements described in Section~\ref{sec:bg:benchmarks} depend on the first property, and are therefore natural for the LLaVA-style architectures used in this thesis. The two backbones evaluated here sit at different points on the remaining spectrum: \texttt{HuatuoGPT-Vision-7B}~\cite{chen2024towards} follows the LLaVA-style projection design, while \texttt{Qwen3-VL-4B-Instruct}~\cite{bai2025qwen3} employs an adapter supporting dynamic, high-resolution inputs, so that the number of visual tokens varies with the input rather than being fixed.

\subsubsection{Autoregressive Generation over Interleaved Sequences}

Once projected, visual tokens are simply prepended to or interleaved with the text token embeddings, and the language model runs unmodified. The model defines a distribution over responses by the usual autoregressive factorization,
\begin{equation}\label{eq:autoregressive}
\pi_\theta(y \mid v, q) = \prod_{t=1}^{|y|} \pi_\theta(y_t \mid v, q, y_{<t}),
\end{equation}
where $v$ denotes the image, $q$ the text prompt, and $y$ the generated response.

Three consequences of Equation~\ref{eq:autoregressive} recur throughout this thesis.

First, the log-likelihood of a response is a \emph{sum} over token positions,
\begin{equation}\label{eq:logprob-sum}
\log \pi_\theta(y \mid v,q) = \sum_{t=1}^{|y|} \log \pi_\theta(y_t \mid v,q,y_{<t}).
\end{equation}
Any objective defined on sequence log-probabilities is therefore implicitly a sum of per-token quantities. This is what makes it possible --- and, as Chapter~\ref{section:analysis} argues, necessary --- to restrict such an objective to a subset of positions: a sum can be taken over an index set, whereas a genuinely sequence-level quantity could not be decomposed at all.

Second, longer responses accumulate more terms, so sequence-level quantities scale with length. A preference objective built on differences of log-probabilities will therefore see systematically larger magnitudes for long responses than short ones, independently of their content. This has direct consequences for gradient magnitudes and is one route by which response length becomes a confound.

Third, the image influences generation only through the attention that text positions pay to the visual token positions. There is no other pathway: once projected, visual tokens are ordinary sequence elements, and the only mechanism by which they can affect the distribution over $y_t$ is the attention operation. This is what makes attention over visual tokens a meaningful proxy for visual grounding, and it is the basis of the metrics described in Section~\ref{sec:bg:benchmarks}. It is also why a model can produce a correct answer with negligible visual attention --- the language model's priors alone may suffice --- which is the failure mode Section~\ref{sec:ana:grounding} is concerned with.

\subsubsection{Instruction Tuning}

A model trained only to predict the next token continues text; it does not answer questions. Instruction tuning closes this gap by supervised fine-tuning on a large collection of (instruction, response) pairs spanning many tasks, which induces the ability to follow instructions expressed in natural language, including for tasks not seen during tuning~\cite{wei2022finetuned}. Visual instruction tuning extends this to the multimodal case by constructing (image, instruction, response) triples, typically by prompting a strong text-only model with image annotations to synthesize plausible dialogues~\cite{liu2023visual, dai2023instructblip}.

Instruction tuning is the last stage that both backbones in this thesis have already undergone, and it is the starting point --- the reference policy $\pi_{\text{ref}}$ --- for everything that follows. Its limitation motivates the rest of this chapter. Supervised fine-tuning teaches the model to imitate the distribution of demonstrated responses, but it provides no mechanism for expressing that one response is \emph{better} than another. Where demonstrations are ambiguous, incomplete, or occasionally wrong, imitation reproduces those properties faithfully.

This limitation is particularly consequential for synthetically generated instruction data, which is how most multimodal instruction corpora are built. When a text-only model writes a dialogue about an image it cannot see, working only from annotations, the resulting response is fluent and well-formed by construction but its grounding in the image is only as good as the annotation. A model trained to imitate such data learns the register of grounded description without necessarily learning the grounding, which is precisely the dissociation between surface form and evidential support that this thesis is about.

\subsection{Medical Vision-Language Models}
\label{sec:bg:medical-lvlm}

\subsubsection{Domain-Specialized Pretraining}

General-purpose LVLMs perform poorly on medical imaging out of the box. Medical images differ from natural images in nearly every respect that matters to a visual encoder. They are frequently grayscale and low in contrast; they are dominated by texture rather than by object boundaries; the diagnostically relevant feature may occupy a small fraction of the frame; and the spatial conventions are rigid in ways natural images are not, so that a left-right flip is not an innocuous augmentation but a change of clinical meaning. The vocabulary describing them is specialized to the point of being nearly disjoint from web caption text, and much of the clinically important content is expressed through negation --- a report stating that a finding is \emph{absent} is making a substantive claim, not declining to make one.

The standard response is continued pretraining on medical image-text corpora. LLaVA-Med~\cite{li2024llavamed} adapts the LLaVA recipe using figure-caption pairs mined from biomedical literature, followed by instruction tuning on synthesized medical dialogues. CheXagent~\cite{chen2024chexagent} specializes to chest radiography with a large curated corpus of radiographs and reports. HuatuoGPT-Vision~\cite{chen2024towards}, one of the two backbones used here, scales this approach with a large collection of medical image-text pairs spanning modalities. MedGemma~\cite{sellergren2025medgemma} applies a comparable strategy to a different base family. Broader surveys document the rapid proliferation of such systems and their proposed clinical applications~\cite{alsaad2024multimodal, saab2024capabilities}.

\subsubsection{What Medical Pretraining Does and Does Not Provide}

Domain pretraining reliably supplies vocabulary, format, and plausibility. It does not reliably supply factuality or grounding, and the gap between these is the problem this thesis addresses.

The most-documented failure is \emph{hallucination}: the model asserts findings that are not present in the image, in fluent and appropriately hedged clinical prose~\cite{zhang2024medihall, seth2025hallucinogen}. Hallucinated content is difficult to detect precisely because domain pretraining has made the surface form correct; the output reads exactly like a report a radiologist would write, and only the content is wrong. Worse, the fluency is actively misleading to a reader, since in human-authored text fluency and confidence correlate with expertise.

Closely related is insensitivity to visual evidence: models frequently produce answers driven by linguistic priors --- the base rate of a finding, the anatomical default for a modality, or a correlation between question phrasing and answer --- rather than by the image~\cite{xia2024cares}. Because a prior-driven answer and an evidence-driven answer are often the same string, this failure is invisible to any evaluation that inspects only outputs. It is also self-concealing in aggregate: a model exploiting base rates will score well on a test set with the same base rates, and will fail specifically on the atypical cases, which are both rare in the test distribution and the ones a clinician most needs help with.

Models also fail to follow clinician instructions reliably, ignoring constraints on scope or format~\cite{vadlapati2026ai, xia2024cares}, and exhibit inconsistency under clinically irrelevant perturbations of the question.

These failures share a structure: they are all cases where the model's output distribution is well-formed but incorrectly conditioned. The model has learned what a good answer looks like without having learned what makes a superficially identical answer wrong. Supervised fine-tuning on more demonstrations does not obviously fix this, because a demonstration exhibits only the correct answer; it never exhibits the near-miss that must be rejected. This is precisely the situation preference-based methods are designed for, since a preference pair carries exactly the information a demonstration lacks --- a negative example that is close enough to the positive one to isolate the distinction.

\subsection{Post-Training Alignment}
\label{sec:bg:alignment}

\subsubsection{Supervised Fine-Tuning}

Supervised fine-tuning (SFT) minimizes the negative log-likelihood of demonstrated responses,
\begin{equation}\label{eq:sft}
\mathcal{L}_{\text{SFT}}(\pi_\theta) = -\mathbb{E}_{(v,q,y) \sim \mathcal{D}}\left[\sum_t \log \pi_\theta(y_t \mid v,q,y_{<t})\right].
\end{equation}

Its limitations are structural rather than practical. It only ever increases the probability of observed responses, so it cannot express that an unobserved response is bad; probability mass is removed from alternatives only incidentally, through normalization. It treats every demonstration as equally authoritative, so noise in the corpus is learned along with signal. And it optimizes an objective --- likelihood of a particular reference --- that is only loosely correlated with the quality a reader cares about, since many acceptable responses receive low likelihood merely by differing in wording from the one that happened to be written down. The alignment literature exists to supply the missing comparative signal.

\subsubsection{Reinforcement Learning from Human Feedback}

Reinforcement learning from human feedback (RLHF) obtains that signal by asking annotators to compare outputs rather than produce them~\cite{christiano2017deep, stiennon2020learning, ouyang2022training}. Comparison is both cheaper and more reliable than demonstration, and the asymmetry is especially pronounced in specialist domains: writing a gold-standard radiology report requires a radiologist, whereas ranking two candidate reports requires less time from the same expert and admits more consistent agreement.

The classical pipeline has two stages. First, a reward model $r_\phi$ is fit to the collected comparisons under the Bradley-Terry model of paired choice~\cite{bradley1952rank}, which posits that the probability of preferring $y^+$ over $y^-$ is a logistic function of their reward difference:
\begin{equation}\label{eq:bt}
p(y^+ \succ y^- \mid v, q) = \frac{\exp r_\phi(v,q,y^+)}{\exp r_\phi(v,q,y^+) + \exp r_\phi(v,q,y^-)} = \sigma\big(r_\phi(v,q,y^+) - r_\phi(v,q,y^-)\big).
\end{equation}
Fitting $r_\phi$ by maximum likelihood over a preference dataset is then ordinary logistic regression on reward differences. Note that the model is identified only up to an additive constant per prompt, since the likelihood depends solely on differences --- a property that becomes essential in the DPO derivation below.

Second, the policy is optimized against this learned reward with a KL penalty anchoring it to the SFT model,
\begin{equation}\label{eq:rlhf}
\max_{\pi_\theta}\; \mathbb{E}_{y \sim \pi_\theta}\big[r_\phi(v,q,y)\big] - \beta\, \mathbb{D}_{\text{KL}}\big(\pi_\theta(y \mid v,q) \parallel \pi_{\text{ref}}(y \mid v,q)\big),
\end{equation}
typically using PPO~\cite{schulman2017proximal}.

The KL term is not decoration. The learned reward $r_\phi$ is accurate only on the distribution of responses the annotators actually saw; far from that distribution it is an extrapolation, and policies are extremely effective at finding regions where an imperfect reward model is erroneously optimistic. Without the penalty the policy drifts to whatever degenerate output maximizes this error, a phenomenon known as reward over-optimization. The penalty confines the search to a neighbourhood of the reference where the reward model's estimates remain trustworthy, and $\beta$ sets the radius of that neighbourhood. The reader should note that the divergence in Equation~\ref{eq:rlhf} is written in the \emph{reverse} direction, $\mathbb{D}_{\text{KL}}(\pi_\theta \parallel \pi_{\text{ref}})$, a choice whose significance is developed in Section~\ref{sec:bg:kl}.

RLHF is effective but operationally heavy: it requires training and serving a separate reward model, sampling from the policy online during training, and absorbing the tuning burden of on-policy reinforcement learning, which is considerably less forgiving than supervised optimization.

\subsubsection{Direct Preference Optimization}

Direct Preference Optimization~\cite{rafailov2023direct} removes the reward model entirely by observing that the optimum of Equation~\ref{eq:rlhf} has a closed form. For any reward function $r$, the KL-constrained maximization is solved by
\begin{equation}\label{eq:optimal-policy}
\pi^{*}(y \mid v,q) = \frac{1}{Z(v,q)}\, \pi_{\text{ref}}(y \mid v,q)\, \exp\!\left(\tfrac{1}{\beta} r(v,q,y)\right),
\end{equation}
where $Z(v,q) = \sum_{y} \pi_{\text{ref}}(y \mid v,q)\exp(r(v,q,y)/\beta)$ normalizes the distribution. The form is intuitive: the optimal policy reweights the reference by an exponential tilt proportional to reward, with $\beta$ controlling how aggressive the tilt is. As $\beta \to \infty$ the policy reverts to the reference; as $\beta \to 0$ it collapses onto the reward maximizer.

Rearranging Equation~\ref{eq:optimal-policy} expresses the reward in terms of the optimal policy:
\begin{equation}\label{eq:reward-from-policy}
r(v,q,y) = \beta \log \frac{\pi^{*}(y \mid v,q)}{\pi_{\text{ref}}(y \mid v,q)} + \beta \log Z(v,q).
\end{equation}

The partition function $Z(v,q)$ is intractable, since it sums over all possible responses. The key step is what happens when Equation~\ref{eq:reward-from-policy} is substituted into the Bradley-Terry likelihood of Equation~\ref{eq:bt}. Because that likelihood depends only on the \emph{difference} of two rewards conditioned on the same $(v,q)$, the intractable term appears in both and cancels exactly. Parameterizing the policy directly as $\pi_\theta$ yields an objective that can be optimized by supervised learning on a static preference dataset:
\begin{equation}\label{eq:dpo-derived}
\mathcal{L}_{\text{DPO}}(\pi_\theta; \pi_{\text{ref}}) = -\mathbb{E}\left[\log \sigma\left(\beta \log \frac{\pi_\theta(y^+ \mid v,q)}{\pi_{\text{ref}}(y^+ \mid v,q)} - \beta \log \frac{\pi_\theta(y^- \mid v,q)}{\pi_{\text{ref}}(y^- \mid v,q)}\right)\right].
\end{equation}

The language model is thus, in the phrase of the original paper, secretly its own reward model: the quantity $\beta \log \pi_\theta / \pi_{\text{ref}}$ is the \emph{implicit reward}.

\subsubsection{The DPO Gradient and What It Implies}

The properties that matter for this thesis are clearest in the gradient. Writing $\hat{r}_\theta(y) = \beta \log \frac{\pi_\theta(y \mid v,q)}{\pi_{\text{ref}}(y \mid v,q)}$ for the implicit reward and $\Delta = \hat{r}_\theta(y^+) - \hat{r}_\theta(y^-)$ for the margin, the gradient of Equation~\ref{eq:dpo-derived} is
\begin{equation}\label{eq:dpo-gradient}
\nabla_\theta \mathcal{L}_{\text{DPO}} = -\beta\,\sigma(-\Delta)\,\Big[\nabla_\theta \log \pi_\theta(y^+ \mid v,q) - \nabla_\theta \log \pi_\theta(y^- \mid v,q)\Big].
\end{equation}

Three observations follow, and each anticipates an argument in Chapter~\ref{section:analysis}.

\textbf{The per-example weight is a scalar that vanishes as the margin grows.} The factor $\sigma(-\Delta)$ is large when the model ranks the pair incorrectly and tends to zero as $\Delta$ becomes large and positive. This is sensible --- already-satisfied examples should stop contributing --- but it also means that a model which finds \emph{any} reliable way to separate the pair will drive its own learning signal to zero. If that separation is a spurious stylistic feature rather than the intended semantic one, training effectively halts with the intended distinction unlearned. Section~\ref{sec:res:dynamics} observes exactly this, with gradient norms collapsing by six orders of magnitude once the margin saturates.

\textbf{The reward is uniform over tokens.} Expanding $\log \pi_\theta(y \mid v,q)$ using Equation~\ref{eq:logprob-sum}, the bracketed term in Equation~\ref{eq:dpo-gradient} is a sum of per-token gradients, every one of which carries the same scalar coefficient $\beta\sigma(-\Delta)$. The objective has no mechanism for distinguishing the tokens that actually differ between $y^+$ and $y^-$ from those the two responses share. When the semantic difference is concentrated in a two-word span inside a thirty-word response, the great majority of the gradient is applied to tokens that carry no information about the preference at all. This is the coarseness that Section~\ref{sec:ana:coarse} identifies and that the fine-grained ranking loss of Section~\ref{sec:met:ranking} addresses.

\textbf{The KL constraint is implicit in the reward's denominator.} Unlike Equation~\ref{eq:rlhf}, the DPO objective contains no explicit divergence term. The constraint survives only through the $\pi_{\text{ref}}$ appearing inside $\hat{r}_\theta$. This has an important consequence for any method that restricts the reward to a subset of token positions: doing so silently removes the KL constraint on all the remaining positions, which is why the fine-grained methods of Section~\ref{sec:bg:dpo-variants} must reintroduce a divergence term explicitly. Section~\ref{sec:res:abl-kl} shows what happens when they do not.

A fourth property is not visible in the gradient but follows from the objective's form: it is a function purely of the \emph{contrast} between $y^+$ and $y^-$. DPO never rewards $y^+$ for being correct in any absolute sense, only for being more likely than $y^-$ under the reweighted policy. Consequently the mechanism by which those two responses were selected determines what the objective actually teaches, which is the subject of Section~\ref{sec:ana:reward-hacking}.

\subsubsection{Extending DPO to the Multimodal Setting}

Applied to an LVLM, Equation~\ref{eq:dpo-derived} simply conditions on $v$ alongside $q$, and nothing in the derivation changes. This is convenient but incomplete: the objective contains no term that references $v$ except through the conditioning, so a policy that ignores the image entirely can still reduce the loss to zero on any preference pair whose members are distinguishable from text alone. Nothing in the mathematics notices the difference between a model that reads the image and one that does not.

mDPO~\cite{wang2024mdpo} addresses this by adding a conditional preference term contrasting the same response under the original and a perturbed image, forcing the model's output distribution to depend on visual input. This introduces two design choices that the original DPO derivation does not constrain: what perturbation to apply, and in which direction the resulting preference should point. Chapter~\ref{section:method} argues that both choices matter considerably in the clinical setting, and Section~\ref{sec:res:abl-visual} measures how much.

\subsection{Beyond Vanilla DPO}
\label{sec:bg:dpo-variants}

\subsubsection{Alternative Preference Objectives}

Several methods retain DPO's reward-model-free structure while changing the loss. IPO~\cite{azar2024general} identifies an overfitting pathology: as the empirical preferences become deterministic, the Bradley-Terry substitution drives the implicit reward difference towards infinity, so the model is pushed to make $\pi_\theta(y^-)$ vanish rather than merely rank it lower. IPO replaces the logistic loss with a bounded squared objective that targets a finite margin. KTO~\cite{ethayarajh2023kto} draws on prospect theory to define a utility over single labelled outputs, removing the requirement for paired data entirely --- valuable when feedback arrives as isolated judgements rather than comparisons.

These address the shape of the objective. The concerns of this thesis are orthogonal to them: they are about which tokens the objective covers and how the pairs are constructed, and would apply equally to an IPO or KTO formulation.

Empirical studies have also documented instability in DPO relative to PPO-based RLHF, particularly the tendency of the policy to diverge and lose general capability~\cite{xu2024dpo, yan20243d}, and have analyzed the specific role of the rejected response in driving that divergence~\cite{cho2025rethinking}. The mechanism is visible in Equation~\ref{eq:dpo-gradient}: nothing bounds how far $\log \pi_\theta(y^-)$ may be driven down, and pushing probability mass off the rejected response redistributes it in ways the objective does not control.

\subsubsection{Token- and Segment-Level Preference Optimization}

A more directly relevant line of work attacks the granularity of the reward. Token-level DPO~\cite{zeng2024token} reformulates preference optimization in the token-level Markov decision process so that credit can be assigned per step rather than per sequence. TGDPO~\cite{zhu2025tgdpo} derives token-level reward guidance to shape the update. Mask-DPO~\cite{gu2025mask} operates at sentence granularity, masking out sentences whose factuality cannot be verified so that unverifiable content does not contribute gradient. ASPO~\cite{wang2025aspo} adaptively weights individual segments according to their estimated contribution.

RRPO~\cite{sarkar2025self}, the most direct antecedent of this thesis, restricts the ranking loss to the tokens that actually differ between the preferred and rejected responses, and compensates for the resulting loss of the implicit KL constraint by adding an explicit forward token-wise KL term. Its formulation is given in full in Section~\ref{sec:ana:prelim}, and the three respects in which it does not transfer to clinical alignment are the subject of Chapter~\ref{section:analysis}.

Common to this family is a dependence on being able to \emph{identify} the tokens that matter. Token-level DPO infers this from the optimization; Mask-DPO uses an external factuality check; RRPO uses a template alignment against a reference. The approach taken in Chapter~\ref{section:method} is to make the identification explicit at data construction time, by generating the preference pair through a minimal edit and recording where that edit occurred.

\subsubsection{On-Policy Preference Construction}

A separate thread concerns not the objective but the data. Because DPO is trained on a fixed dataset, the responses in that dataset are generally not samples from the policy being trained, and this off-policy mismatch degrades alignment quality~\cite{yang2025mitigating}. The concern is partly statistical --- the derivation assumes preferences over the policy's own distribution --- and partly practical, in that a preferred response the model would essentially never generate provides a target it cannot reach by small parameter changes.

OPA-DPO~\cite{yang2025opadpo} argues for constructing preferred responses by sampling from the model and minimally revising them, so that the target remains within the model's reachable manifold. This thesis adopts the same principle but arrives at it from a different direction. Section~\ref{sec:ana:reward-hacking} shows that in the clinical setting the mismatch takes a specific and damaging form: a systematic stylistic gap between clinical references and model outputs, which the model exploits in preference to learning the clinical distinction. On-policy construction is therefore not merely a matter of reachability but of removing a confound that would otherwise dominate the learning signal --- a claim measured directly in Section~\ref{sec:res:style-gap}.

\subsection{Information-Theoretic Preliminaries}
\label{sec:bg:kl}

The Kullback-Leibler divergence between distributions $P$ and $Q$ over a discrete space is
\begin{equation}\label{eq:kl-def}
\mathbb{D}_{\text{KL}}(P \parallel Q) = \sum_{x} P(x) \log \frac{P(x)}{Q(x)},
\end{equation}
which is non-negative and zero if and only if $P = Q$~\cite{kullback1951information}. It is not symmetric, and the asymmetry is not a technicality: the two orderings induce qualitatively different optimal solutions whenever the approximating family cannot represent the target exactly~\cite{minka2005divergence}. Since a LoRA-adapted policy has far fewer effective degrees of freedom than the space of distributions over token sequences, this regime is exactly the one alignment operates in.

\subsubsection{Forward KL and Mean-Seeking Behaviour}

Taking $P = \pi_{\text{ref}}$ and $Q = \pi_\theta$ gives the \emph{forward} direction, $\mathbb{D}_{\text{KL}}(\pi_{\text{ref}} \parallel \pi_\theta)$. The expectation in Equation~\ref{eq:kl-def} is taken under the reference, so the divergence is only evaluated where the reference has mass. Wherever $\pi_{\text{ref}}(x) > 0$ and $\pi_\theta(x) \to 0$, the log ratio diverges and the penalty is unbounded; but wherever $\pi_\theta(x) > 0$ and $\pi_{\text{ref}}(x) = 0$, the term is weighted by $P(x) = 0$ and contributes nothing at all.

Forward KL therefore forbids \emph{dropping} mass and is indifferent to \emph{adding} it. The optimum is \emph{mean-seeking} or \emph{mass-covering}: when the approximating family is too restrictive to match a multimodal target, the fitted distribution stretches to cover every mode, necessarily placing mass in the low-density regions between them. Discarding the terms that do not depend on $\theta$ reduces the objective to maximum likelihood under samples from the reference, which is the same functional form as Equation~\ref{eq:sft} --- forward KL is, in effect, distillation of the reference into the policy.

\subsubsection{Reverse KL and Mode-Seeking Behaviour}

Reversing the arguments gives $\mathbb{D}_{\text{KL}}(\pi_\theta \parallel \pi_{\text{ref}})$, with the expectation now taken under the policy. The asymmetry inverts: the policy is heavily penalized for placing mass where the reference assigns little, and pays nothing for ignoring regions the reference covers. The optimum is \emph{mode-seeking} or \emph{zero-forcing}, concentrating on one high-density region and abandoning the rest.

Expanding the definition gives
\begin{equation}\label{eq:rkl-decomp}
\mathbb{D}_{\text{KL}}(\pi_\theta \parallel \pi_{\text{ref}}) = -\mathbb{E}_{y \sim \pi_\theta}\big[\log \pi_{\text{ref}}(y)\big] - \mathcal{H}(\pi_\theta),
\end{equation}
where $\mathcal{H}$ denotes entropy. The decomposition separates two competing pressures: the first term rewards the policy for placing mass where the reference finds it likely, and the second, the entropy, opposes total collapse onto a single output. Which dominates determines whether the result is a sharpened policy or a degenerate one.

The expectation being taken under $\pi_\theta$ is what makes reverse KL an \emph{on-policy} regularizer: it evaluates the reference only at points the current policy actually visits, so the penalty adapts as the policy moves. This is also why it is the direction used in the RLHF objective of Equation~\ref{eq:rlhf} and, implicitly, in DPO.

\subsubsection{Implications for Factuality-Critical Generation}

The distinction is consequential when the objective is factual correctness rather than general quality.

The characteristic failure of a medical LVLM is a fluent, confident, clinically wrong continuation. In distributional terms, such a continuation is a token sequence to which the reference model assigns low but non-negligible probability --- it is not absurd, or it would never be generated, but it is not the mode either. Suppressing exactly this kind of tail is what reverse KL does and what forward KL structurally cannot: a forward penalty is indifferent to whether the policy retains, or even amplifies, a mode that the reference finds unlikely, because such regions carry weight zero in its expectation.

The converse also holds, and is equally important. A purely mode-seeking regularizer will narrow the policy's output distribution, which matters when the same model must handle open-ended questions that admit several correct phrasings and require varied clinical reasoning. Collapse also interacts badly with the preference objective itself: a policy with very low entropy produces near-identical samples, so any subsequent on-policy data construction degenerates.

Neither direction is adequate alone. Chapter~\ref{section:analysis} makes this trade-off empirical, and Section~\ref{sec:met:kl} develops the bidirectional regularizer that resolves it by interpolating between the two with a single coefficient.

\subsection{Medical Benchmarks and Evaluation}
\label{sec:bg:benchmarks}

\subsubsection{Medical Visual Question Answering}

Medical VQA presents a model with a clinical image and a natural-language question. Two datasets are standard in this literature.

VQA-RAD~\cite{lau2018dataset} contains radiology images of head, chest, and abdomen paired with question-answer pairs written by clinicians. It is valued for the authenticity of its questions --- they reflect what a physician would actually ask of an image --- rather than for its scale, which is modest.

SLAKE~\cite{liu2021slake} is larger and bilingual, and is additionally annotated with segmentation masks and a medical knowledge graph. This thesis uses only its English visual-textual pairs; the knowledge graph is deliberately excluded, since it would supply structured clinical information that the preference construction pipeline of Chapter~\ref{section:method} is meant to derive from the image and the ground-truth answer alone.

Both datasets distinguish \emph{closed} questions, with a binary or categorical answer, from \emph{open} questions requiring free-form generation. The distinction matters for evaluation in a way that is easy to underestimate. A closed question admits a small answer set and is therefore partially answerable by prior: a model that always answers ``no'' to presence questions will score at the base rate without reading the image at all. Open questions require the model to produce the clinical content itself and are correspondingly harder to satisfy by guessing, which is why this thesis reports the two separately throughout and treats improvements on open questions as the stronger evidence.

Broader benchmarks such as PathVQA~\cite{he2020pathvqa} and PMC-VQA~\cite{zhang2023pmcvqa} extend the setting to other modalities and are natural candidates for evaluating whether alignment gains transfer beyond the corpus a model was tuned on.

\subsubsection{Radiology Report Generation}

Report generation asks for a full free-text radiology report given one or more views. IU-Xray~\cite{demner2016preparing} pairs dual-view chest radiographs with their corresponding reports and is the standard benchmark for the task.

The task differs from VQA in a way that matters throughout this thesis. Outputs are long and contain many independent clinical assertions, any one of which may be wrong while the rest are correct. A single scalar preference over two such reports is a very coarse instrument: it compresses a dozen independent judgements into one bit. Report generation is therefore the regime where the granularity problem of Section~\ref{sec:ana:coarse} is most acute, and, as Section~\ref{sec:res:main} shows, the regime where the methods of this thesis help most.

\subsubsection{Evaluation Metrics}

Free-form clinical text cannot be scored by exact match. The conventional alternative, n-gram overlap metrics such as BLEU~\cite{papineni2002bleu} and ROUGE~\cite{lin2004rouge}, is actively misleading here: inserting a negation changes almost no tokens while inverting the clinical meaning, so a report can score highly while being diagnostically the opposite of correct. Since negation carries much of the clinical content in radiology reporting, this is not an edge case.

Two alternatives are used in this thesis.

For VQA, a strong language model serves as a judge, presented with the question, the ground-truth answer, and the model's response, and asked to decide correctness. This accommodates paraphrase --- ``left lung'' and ``the lesion is in the left lung'' are the same answer --- while remaining sensitive to clinical content in a way string matching is not.

For report generation, the GREEN score~\cite{ostmeier2024green} uses a model fine-tuned to identify and categorize clinically significant errors between a candidate and a reference report, yielding accuracy, precision, and recall over clinical findings rather than over n-grams. Because it operates on findings, it is insensitive to the phrasing differences that dominate n-gram metrics and sensitive to exactly the substitutions that matter.

Both metrics are themselves model-based, which is a genuine limitation rather than an incidental one. Language-model judges are known to disagree with clinicians in systematic ways~\cite{delucia2026same}, and a judge that shares architecture or training data with the model under evaluation may share its blind spots. The results in Chapter~\ref{section:results} should be read with that caveat in mind.

\subsubsection{Measuring Visual Grounding}

Output-level metrics cannot distinguish an answer derived from the image from an identical answer derived from a prior. Since that distinction is central to this thesis, grounding must be measured inside the model rather than at its output.

The approach used here exploits the observation from Section~\ref{sec:bg:lvlm} that the image influences generation only through attention to visual token positions. If a text token's prediction depends on the image at all, that dependence is mediated by attention weights over the visual positions, and those weights are observable.

The \emph{Attention Ratio}~\cite{khayatkhoei2025mllms} quantifies what fraction of a text token's attention over visual positions falls inside the region relevant to the question, normalized by what a uniform distribution would give. A value of one indicates attention no better than chance; values above one indicate concentration on the relevant region. Comparing the model's attention distribution against the reference region distribution via KL divergence, and via the symmetric Jensen-Shannon divergence, provides two further measures, where lower values indicate closer agreement.

VGMED~\cite{liu2026vgmed} operationalizes this protocol as a benchmark built by its authors over a subset of SLAKE, supplying the per-question region annotations the measurement requires and separating clinically critical tokens into Localization and Attribute categories. All three quantities are computed per decoder layer, which allows the depth at which grounding emerges to be examined rather than only its final value --- a distinction that turns out to matter in Section~\ref{sec:res:grounding}, where the methods separate most clearly in the deeper layers.
